\section{Discussion}
\label{sec:discussion}

\paragraph{What made directing work.} Three choices mattered most. First, targets came from code \emph{and} history:
the bugs of Method D sit next to recent fixes (\#28541 is the UHUGEINT sibling of the UBIGINT bug \#28535, found one day later by the loop on the same rule). Second, the LLM had to state a hypothesis about a value the rule builds, which
pushed it toward type limits, NaN and signed zero instead of generic queries. Third, the database judged the result.
Without the judge, the LLM's own reports in Method C would have counted trivial generators as thorough.

\paragraph{Why keep generic fuzzing?} Method A found no new bug, but it is the only method that looks for memory
errors under sanitizers, and its results define the baseline: they show that the bugs of Methods B to D are not ones a
generic fuzzer finds in the same time.

\paragraph{Why not report every difference?} Most differences are not bugs. Of 138 triaged hits in \texttt{hunt1},
only 9 were new candidates. Many of the rest are expected: when the unoptimized plan fails and the optimized plan
succeeds, the optimizer only removed work whose error does not matter, which DuckDB accepts. Sending such reports would
cost maintainers' time and our credibility. We also grouped bugs with one root cause into one issue and filed at a
steady pace.

\paragraph{Limits of the bug family.} Most of our bugs belong to one family: constants at type limits in optimizer
rewrites. Yield per target falls once a rule's family is exhausted (Table~\ref{tab:hunt1}). New families need new
oracles: TLP and NoREC can find bugs that are present with the switch both on and off, and a cross-version oracle
(release against development build) finds regressions directly. Rules that match operators such as \texttt{+} or
\texttt{LIKE} have no marker function, so their reach is measured only at expression level. Measuring the coverage of
the target file's lines would give a stronger signal for them. We leave these for future work.
